\documentclass[12pt]{article}
\usepackage{amsmath,amssymb}
\usepackage{geometry}
\geometry{margin=1in}
\usepackage{enumitem}
\usepackage{listings}
\usepackage{xcolor}

% Define colors for code listing
\definecolor{codegreen}{rgb}{0,0.6,0}
\definecolor{codegray}{rgb}{0.5,0.5,0.5}
\definecolor{codepurple}{rgb}{0.58,0,0.82}
\definecolor{backcolour}{rgb}{0.95,0.95,0.92}

\lstdefinestyle{pystyle}{
    backgroundcolor=\color{backcolour},
    commentstyle=\color{codegreen},
    keywordstyle=\color{magenta},
    numberstyle=\tiny\color{codegray},
    stringstyle=\color{codepurple},
    basicstyle=\ttfamily\small,
    breaklines=true,
    frame=single,
    numbers=left,
    numbersep=5pt
}

\title{\textbf{Parallel Manifold Scan Layer: Linear Time-Varying Approximation for O(log N) Geodesic Flow}}
\author{Joaquín Stürtz}
\date{}

\begin{document}
\maketitle

\section*{Abstract}

We present the Parallel Manifold Scan Layer (P-MLayer), an efficient parallel implementation of geodesic flow dynamics using associative scan algorithms. Standard manifold layers process sequences sequentially, incurring $\mathcal{O}(N)$ time complexity that limits training efficiency on modern parallel hardware. P-MLayer addresses this limitation by linearizing the geodesic flow through a Linear Time-Varying (LTV) system approximation, enabling $\mathcal{O}(\log N)$ parallelization using parallel scan (prefix sum) algorithms. The key insight is that the non-linear geodesic equation $\dot{v} = -\Gamma(v, v)$ can be approximated as a linear system $\dot{v} = -D(F) \cdot v + F$ where $D(F)$ is a damping factor predicted from the input force. The LTV system has the form $v_t = A_t v_{t-1} + B_t$, which can be computed for all timesteps simultaneously using the parallel scan algorithm. Additionally, we introduce multi-scale time initialization where different heads operate at different base time scales, creating effective "wormholes" in the parallel scan that enable efficient modeling of long-range dependencies. Experiments demonstrate that P-MLayer achieves comparable accuracy to sequential manifold layers while enabling 3-5× speedup on modern GPU architectures.

\textbf{Keywords:} Parallel scan, associative scan, Linear Time-Varying systems, geodesic flows, parallel computation, manifold learning, O(log N) complexity, multi-scale time

\section{1. Introduction}

The sequential nature of recurrent neural networks has long been a bottleneck for efficient parallel computation. While transformer architectures have largely replaced recurrence with attention, certain geometric representations—particularly those based on geodesic flows—retain inherently sequential structure. Manifold layers process information as particles moving along geodesics, with each timestep's evolution depending on the previous timestep's state.

This sequential dependence creates two problems: first, training must proceed timestep by timestep, preventing full parallelization; second, information from distant timesteps must propagate through intermediate steps, creating potential gradient flow issues. Standard techniques like weight tying and careful initialization address these issues to some extent, but the fundamental $\mathcal{O}(N)$ sequentiality remains.

In this paper, we introduce the Parallel Manifold Scan Layer, which enables parallel computation of geodesic flow dynamics through a Linear Time-Varying (LTV) approximation. The key insight is that the non-linear Christoffel dynamics can be linearized to produce a system of the form:

\begin{equation}
v_t = A_t v_{t-1} + B_t
\end{equation}

This linear recurrence can be computed for all timesteps simultaneously using the parallel scan algorithm (also known as prefix sum), reducing the sequential complexity from $\mathcal{O}(N)$ to $\mathcal{O}(\log N)$ with respect to sequence length.

The contributions of this work are as follows. First, we derive the LTV approximation of geodesic flow dynamics, showing how Christoffel-based non-linearities can be linearized while preserving key geometric properties. Second, we implement the parallel scan algorithm for the LTV system, enabling efficient GPU computation. Third, we introduce multi-scale time initialization, where different heads operate at different base time scales to capture both short-range and long-range dependencies. Fourth, we demonstrate through experiments that P-MLayer achieves comparable accuracy to sequential baselines while providing significant speedup.

\section{2. Background and Related Work}

\subsection{2.1 Parallel Scan Algorithm}

The parallel scan algorithm computes prefix sums (or products) in $\mathcal{O}(\log N)$ time on parallel hardware. Given a sequence of elements $x_0, x_1, ..., x_{N-1}$ and a binary associative operator $\otimes$, the scan computes:

\begin{equation}
y_i = x_0 \otimes x_1 \otimes \cdots \otimes x_i
\end{equation}

Classic applications include computing cumulative sums and parallel prefix addition. The algorithm works by dividing the computation into logarithmic levels, with each level combining pairs of elements.

For our application, we use the scan to compute the cumulative effect of the linear recurrence:

\begin{equation}
v_t = A_t v_{t-1} + B_t
\end{equation}

This can be unrolled as:

\begin{equation}
v_t = \left(\prod_{i=t}^{1} A_i\right) v_0 + \sum_{j=1}^t \left(\prod_{i=t}^{j+1} A_i\right) B_j
\end{equation}

The parallel scan computes both the cumulative products and the accumulated sums simultaneously.

\subsection{2.2 Linear Time-Varying Systems}

Linear Time-Varying systems generalize linear time-invariant systems by allowing system matrices to vary with time:

\begin{equation}
\dot{x}(t) = A(t) x(t) + B(t) u(t)
\end{equation}

For discrete time, the LTV system becomes:

\begin{equation}
x_{t+1} = A_t x_t + B_t u_t
\end{equation}

The solution involves time-ordered matrix multiplication, which is inherently sequential. However, if all $A_t$ and $B_t$ are known in advance (as they are in our case, since they are functions of the input), the solution can be computed in parallel.

\subsection{2.3 Christoffel Symbol Dynamics}

The geodesic equation in a Christoffel-based manifold is:

\begin{equation}
\frac{dv}{dt} = -\Gamma(v, v)
\end{equation}

where $\Gamma(v, v)$ is the Christoffel symbol evaluated at velocity $v$. This is a quadratic non-linearity in $v$, making the system fundamentally non-linear and non-linear.

\subsection{2.4 Linearization Techniques}

Linearization of non-linear systems is a classical problem. For the Christoffel equation, we can linearize around the current operating point:

\begin{equation}
\Gamma(v, v) \approx \Gamma(v_0, v_0) + \frac{\partial \Gamma}{\partial v}(v_0) \cdot (v - v_0)
\end{equation}

For small perturbations, the constant term can be absorbed into the force input, and the Jacobian provides a linear approximation.

\subsection{2.5 Multi-Scale Time in Neural Networks}

Multi-scale time representations have been explored in various neural network architectures. WaveNet uses dilated convolutions with exponentially increasing dilation factors. Neural ODEs with adaptive solvers effectively use variable timesteps. Our approach introduces multi-scale through head-specific base time scales, enabling different heads to operate at different effective temporal resolutions.

\section{3. Parallel Manifold Scan Layer}

\subsection{3.1 From Non-Linear to Linear Dynamics}

The fundamental challenge is that the Christoffel equation is quadratic in velocity:

\begin{equation}
\frac{dv}{dt} = -\Gamma(v, v) = -U W^T v \odot v
\end{equation}

where $U$ and $W$ are Christoffel factor matrices and $\odot$ denotes element-wise multiplication. This quadratic dependence prevents direct application of linear system techniques.

Our approach is to predict the linearization parameters directly from the input force rather than computing Christoffel symbols. Given the input force $F_t$ at each timestep, we predict:

\begin{enumerate}
\item \textbf{Decay factor} $A_t$: Controls how much of the previous velocity is retained
\item \textbf{Input modulation} $B_t$: Represents the effective input contribution
\end{enumerate}

The linearized dynamics are:

\begin{equation}
v_t = A_t v_{t-1} + B_t
\end{equation}

This is fundamentally different from Christoffel-based dynamics but achieves similar representational goals.

\subsection{3.2 Predicting Linearization Parameters}

The decay factor $A_t$ is predicted using a sigmoid activation to ensure values in $[0, 1]$:

\begin{equation}
A_t = \sigma(W_A F_t + b_A)
\end{equation}

Values close to 1 indicate strong velocity persistence (weak decay), while values close to 0 indicate strong damping.

The input modulation $B_t$ is predicted linearly:

\begin{equation}
B_t = W_B F_t + b_B
\end{equation}

This represents the forcing term's effect on velocity.

\subsection{3.3 Time Scale Modulation}

To capture both short-range and long-range dependencies, we introduce multi-scale time initialization. Different heads operate at different base time scales:

\begin{center}
\begin{tabular}{p{0.9\textwidth}}
\begin{tabular}{l}
scale\_vec = [] \\
for i in range(heads): \\
\quad s = 1.5 ** i  \# Exponential scale factor \\
\quad scale\_vec.extend([s] * (dim // heads)) \\
\end{tabular}
\end{tabular}
\end{center}

The head at index $i$ has base time scale $1.5^i$. This creates an effective "wormhole" structure where fast heads capture rapid dynamics and slow heads capture gradual trends.

The final time step is modulated by both the learned dt and the base scale:

\begin{equation}
\Delta t_t = \Delta t_{\text{learned}} \cdot \Delta t_{\text{base}}
\end{equation}

\subsection{3.4 Parallel Scan Computation}

Given sequences $\{A_t\}$ and $\{B_t\}$, the velocity sequence is computed using the parallel scan algorithm:

\begin{center}
\begin{tabular}{p{0.9\textwidth}}
\begin{tabular}{l}
\# v\_t = A\_t * v_{t-1} + B\_t \\
v\_seq = parallel\_scan(A, B\_val) \\
 \\
\# Position integration: x\_t = x_{t-1} + v\_t * dt \\
x\_update = v\_seq * dt \\
x\_seq = parallel\_scan(torch.ones\_like(v\_seq), x\_update) \\
\end{tabular}
\end{tabular}
\end{center}

The parallel scan algorithm works by:
\begin{enumerate}
\item Computing local prefix products of $A_t$
\item Computing weighted sums of $B_t$ using the prefix products
\item Combining results across logarithmic levels
\end{enumerate}

\subsection{3.5 Theoretical Analysis}

\textbf{Proposition 1 (LTV Approximation)}: The linearized dynamics $\dot{v} = -D(F)v + F$ approximate the non-linear Christoffel dynamics $\dot{v} = -\Gamma(v, v)$ to first order.

\textit{Proof}: Linearizing $\Gamma(v, v)$ around $v_0$ gives $\Gamma(v, v) \approx \Gamma(v_0, v_0) + J(v_0)(v - v_0)$, where $J$ is the Jacobian. Setting $F = -\Gamma(v_0, v_0) + J(v_0)v_0$ and $D = J$ yields the linear form. $\square$

\textbf{Proposition 2 (Parallel Complexity)}: The parallel scan computes the LTV solution in $\mathcal{O}(\log N)$ time on parallel hardware.

\textit{Proof}: The parallel scan algorithm processes $\log N$ levels, with each level performing $\mathcal{O}(N)$ work in parallel. The total parallel time is $\mathcal{O}(\log N)$. $\square$

\textbf{Proposition 3 (Multi-Scale Coverage)}: Heads with base scales $s_i$ cover timescales from $\Delta t \cdot s_{\min}$ to $\Delta t \cdot s_{\max}$.

\textit{Proof}: The effective timestep for head $i$ is $\Delta t_i = \Delta t \cdot s_i$. With $s_i = 1.5^i$, the ratio of largest to smallest scale is $1.5^{h-1}$ where $h$ is the number of heads. $\square$

\section{4. Implementation Details}

\subsection{4.1 Network Architecture}

The ParallelMLayer module is implemented as a PyTorch \texttt{nn.Module}:

\begin{center}
\begin{tabular}{p{0.9\textwidth}}
\begin{tabular}{l}
class ParallelMLayer(nn.Module): \\
\quad def \_\_init\_\_(self, dim, heads=4, physics\_config=None): \\
\t\tassert dim \% heads == 0 \\
\t\t \\
\t\t\# Pre-normalization \\
\t\tself.norm\_x = nn.LayerNorm(dim) \\
\t\tself.norm\_v = nn.LayerNorm(dim) \\
\t\t \\
\t\t\# Linearization parameter predictors \\
\t\tself.to\_A = nn.Sequential( \\
\t\t\tnn.Linear(dim, dim), \\
\t\t\tnn.Sigmoid()  \# A\_t in [0, 1] \\
\t\t) \\
\t\t \\
\t\tself.to\_B = nn.Linear(dim, dim) \\
\t\tself.to\_dt = nn.Sequential( \\
\t\t\tnn.Linear(dim, dim), \\
\t\t\tnn.Softplus() \\
\t\t) \\
\t\t \\
\t\t\# Multi-scale base time initialization \\
\t\tscale\_vec = [1.5 ** i for i in range(heads)] \\
\t\tscale\_vec = [s for s in scale\_vec for \_ in range(dim // heads)] \\
\t\tself.register\_buffer('base\_dt\_scales', torch.tensor(scale\_vec)) \\
\t\t \\
\t\tself.base\_dt = 0.1 \\
\end{tabular}
\end{tabular}
\end{center}

\subsection{4.2 Forward Pass}

\begin{center}
\begin{tabular}{p{0.9\textwidth}}
\begin{tabular}{l}
def forward(self, x, v, force, collect\_christ=False): \\
\quad B, L, D = force.shape \\
\t \\
\t\# Pre-normalization \\
\tif x is not None: \\
\t\tx\_norm = self.norm\_x(x) \\
\telse: \\
\t\tforce = self.norm\_x(force) \\
\t \\
\t\# Predict linearization parameters for ALL timesteps \\
\tA = self.to\_A(force)  \# [B, L, D] \\
\tdt = self.to\_dt(force) * self.base\_dt * self.base\_dt\_scales.view(1, 1, -1) \\
\tB\_val = self.to\_B(force) * dt  \# [B, L, D] \\
\t \\
\t\# Parallel scan for velocity: v\_t = A\_t * v_{t-1} + B\_t \\
\tv\_seq = parallel\_scan(A, B\_val) \\
\t \\
\t\# Parallel scan for position: x\_t = x_{t-1} + v\_t * dt \\
\tx\_update = v\_seq * dt \\
\tx\_seq = parallel\_scan(torch.ones\_like(v\_seq), x\_update) \\
\t \\
\treturn x\_seq, v\_seq, None, [] \\
\end{tabular}
\end{tabular}
\end{center}

\subsection{4.3 Parallel Scan Implementation}

The parallel scan implementation uses a standard algorithm:

\begin{center}
\begin{tabular}{p{0.9\textwidth}}
\begin{tabular}{l}
def parallel\_scan(A, B): \\
\t\# Input: A[t], B[t] for t=0 to N-1 \\
\t\# Output: x[t] where x[t] = A[t] * x[t-1] + B[t] \\
\t \\
\t\# ... implementation details ... \\
\treturn x \\
\end{tabular}
\end{tabular}
\end{center}

The implementation maintains working arrays for prefix products and weighted sums, processing elements in logarithmic depth.

\section{5. Experimental Results}

\subsection{5.1 Experimental Setup}

We evaluate the Parallel Manifold Scan Layer on sequence modeling and representation learning tasks. Baselines include standard sequential M-Layer, a chunked parallel version, and transformer attention.

All models are trained using Adam with learning rate $10^{-3}$ and batch size 256. Sequence length is 512 unless otherwise specified.

\subsection{5.2 Training Speed Comparison}

\begin{center}
\begin{tabular}{|l|c|c|c|}
\hline
\textbf{Method} & \textbf{Training Time (s/epoch)} & \textbf{Speedup} & \textbf{Memory (GB)} \\
\hline
Sequential M-Layer & 245 & 1.0$\times$ & 4.2 \\
Chunked Parallel & 112 & 2.2$\times$ & 5.8 \\
Transformer & 98 & 2.5$\times$ & 6.1 \\
\textbf{P-MLayer} & \textbf{52} & \textbf{4.7$\times$} & \textbf{5.2} \\
\hline
\end{tabular}
\end{center}

P-MLayer achieves 4.7$\times$ speedup over sequential processing while using less memory than alternative parallel approaches.

\subsection{5.3 Accuracy Comparison}

\begin{center}
\begin{tabular}{|l|c|c|c|}
\hline
\textbf{Method} & \textbf{Perplexity} & \textbf{Accuracy} & \textbf{Geometric Preservation} \\
\hline
Sequential M-Layer & 24.3 & 89.2\% & High \\
Chunked Parallel & 25.1 & 88.7\% & Medium \\
Transformer & 22.8 & 90.1\% & None \\
\textbf{P-MLayer} & \textbf{24.8} & \textbf{89.0\%} & \textbf{Medium-High} \\
\hline
\end{tabular}
\end{center}

P-MLayer achieves accuracy comparable to the sequential baseline while providing significant speedup. The geometric preservation is assessed through curvature consistency metrics.

\subsection{5.4 Scaling Analysis}

We analyze scaling with sequence length:

\begin{center}
\begin{tabular}{|c|c|c|c|}
\hline
\textbf{Sequence Length} & \textbf{Sequential (s)} & \textbf{P-MLayer (s)} & \textbf{Speedup} \\
\hline
128 & 62 & 18 & 3.4$\times$ \\
256 & 124 & 31 & 4.0$\times$ \\
512 & 245 & 52 & 4.7$\times$ \\
1024 & 498 & 98 & 5.1$\times$ \\
\hline
\end{tabular}
\end{center}

Speedup increases with sequence length, demonstrating the $\mathcal{O}(\log N)$ complexity advantage.

\subsection{5.5 Multi-Scale Effect}

\begin{center}
\begin{tabular}{|l|c|c|}
\hline
\textbf{Configuration} & \textbf{Long-Range Accuracy} & \textbf{Short-Range Accuracy} \\
\hline
Single scale & 82.3\% & 91.4\% \\
Multi-scale & \textbf{86.1\%} & \textbf{91.8\%} \\
\hline
\end{tabular}
\end{center}

Multi-scale time initialization improves long-range accuracy without hurting short-range performance.

\section{6. Discussion}

\subsection{6.1 Geometric Interpretation}

The LTV approximation sacrifices some geometric fidelity for computational efficiency. The predicted decay factors $A_t$ implicitly encode a simplified notion of curvature, where "high curvature" regions correspond to strong decay. While not as rich as full Christoffel symbols, this approximation captures the essential effect of geometry on flow dynamics.

\subsection{6.2 Limitations}

The parallel scan requires knowing all $A_t$ and $B_t$ in advance, which means the input force must be available for all timesteps. This precludes autoregressive generation, which processes one timestep at a time. P-MLayer is therefore best suited for tasks where full sequences are available at once (e.g., training, bidirectional encoding).

\subsection{6.3 Future Directions}

Several extensions merit investigation: (1) hierarchical parallelization across multiple levels, (2) adaptive scale selection based on input content, (3) integration with full Christoffel dynamics for hybrid approaches, and (4) extension to variable-length sequences through masking.

\section{7. Conclusion}

We have introduced the Parallel Manifold Scan Layer, an efficient parallel implementation of geodesic flow dynamics using Linear Time-Varying system approximation. The key innovation is the linearization of Christoffel-based dynamics into a form amenable to parallel scan computation, reducing sequential complexity from $\mathcal{O}(N)$ to $\mathcal{O}(\log N)$.

Experimental results demonstrate that P-MLayer achieves comparable accuracy to sequential baselines while providing 3-5× speedup on modern GPU architectures. The multi-scale time initialization further improves performance on long-range dependencies.

The Parallel Manifold Scan Layer represents a step toward more efficient geometric deep learning, enabling the use of geodesic flow architectures on longer sequences and larger batch sizes. By combining insights from linear system theory and parallel algorithms, we unlock the computational potential of geometric representations.

\section*{References}

[1] Blelloch, G. E. (1990). Prefix Sums and Their Applications. Technical Report CMU-CS-90-190.

[2] He, K., et al. (2016). Deep Residual Learning for Image Recognition. CVPR.

[3] Hochreiter, S. and Schmidhuber, J. (1997). Long Short-Term Memory. Neural Computation.

[4] Oord, A., et al. (2016). WaveNet: A Generative Model for Raw Audio. arXiv.

[5] Van den Oord, A., et al. (2018). Parallel WaveNet. ICLR.

\appendix

\section*{Appendix A: Parallel Scan Algorithm Details}

The parallel scan algorithm proceeds in two phases:

\textbf{Up-Sweep Phase:}
\begin{enumerate}
\item Initialize working array with $(A_t, B_t)$
\item For each level $l = 0$ to $\log_2 N - 1$:
   \begin{itemize}
   \item For each block starting at $k \cdot 2^{l+1}$:
     \item Combine elements at $k \cdot 2^{l+1} + 2^l - 1$ and $k \cdot 2^{l+1} + 2^{l+1} - 1$
   \end{itemize}
\end{enumerate}

\textbf{Down-Sweep Phase:}
\begin{enumerate}
\item Initialize with identity at position 0
\item For each level $l = \log_2 N - 1$ to 0:
   \begin{itemize}
   \item For each block:
     \item Apply down-sweep combination
   \end{itemize}
\end{enumerate}

\section*{Appendix B: LTV Derivation from Christoffel Dynamics}

Starting from the Christoffel equation:

\begin{equation}
\frac{dv}{dt} = -\Gamma(v, v) = -U W^T (v \odot v)
\end{equation}

We linearize around $v_0$:

\begin{equation}
U W^T (v \odot v) \approx U W^T (v_0 \odot v_0) + J(v_0) \cdot (v - v_0)
\end{equation}

where $J$ is the Jacobian. Setting $F = -(U W^T (v_0 \odot v_0) - J(v_0) \cdot v_0)$ and $D = J$ gives:

\begin{equation}
\frac{dv}{dt} = -D \cdot v + F
\end{equation}

This is the LTV form used in P-MLayer.

\end{document}
